\subsection{Mémoire accumulée et résolvantes du transport}
\label{subsec:memoria-resolvente}

Le retour du calendrier n'élimine pas les incréments produits pendant
le parcours. Pour exprimer cette différence en termes opératoriels, on construit
d'abord un registre cumulatif des événements de la trajectoire, puis sa
série génératrice. La réduction des variables de mémoire s'effectuera sur
ce transport, en conservant également la source et le lecteur de sortie.

\subsubsection{Événements orientés et actualisation du registre}

Soit \((d_t)\) la suite des codes entiers de direction conservée
par une trajectoire du TPK. Le code \(d_t\) et la direction cardinale
du curseur sont des coordonnées distinctes ; aucune identification
entre leurs alphabets n'est supposée. Numérotons à partir de zéro les fenêtres
\[
 I_m=\{9m+1,\ldots,9m+9\},\qquad m\geq0.
\]
Dans le premier cycle, \(I_m\) est la fenêtre \(E_{m+1}\) du registre
dodécaphasique. Son orientation est
\(\sigma=(1,1,1,-1,-1,-1,1,1,1,-1,-1,-1)\).
Pour une histoire prolongée, \(\sigma_m\in\{-1,1\}\) est le signe
conservé dans la fenêtre correspondante. L'événement signé est
\begin{equation}
 a_m=\sigma_m\sum_{t\in I_m}
       \mathbf1_{\{2,4,6,8\}}(d_t),\qquad |a_m|\leq9.
 \label{mem:evento}
\end{equation}
Chaque terme est incorporé lorsque sa transition se produit. Cette définition
compte les événements de la trace ; elle ne choisit pas les directions à partir d'une valeur
à obtenir ultérieurement.

Soient \(\mathbf e_0,\ldots,\mathbf e_{11}\) la base de \(\ZZ^{12}\)
et \(S\mathbf e_j=\mathbf e_{j+1\bmod12}\). Écrivons
\(R=S^{-1}\). Le registre accumulé en coordonnées fixes satisfait
\[
z_{m+1}=z_m+a_m\mathbf e_{m\bmod12}.
\]
Un préfixe sans histoire accumulée commence en \(z_0=0\) ; à une autre
origine, \(z_0\) conserve le registre antérieur.
Le repère qui accompagne le calendrier est \(y_m=S^{-m}z_m\).

\begin{proposition}[Actualisation et retour avec mémoire]
\label{mem:prop-actualizacion}
Le registre dans le repère mobile satisfait
\begin{equation}
 y_{m+1}=Ry_m+\eta_m,\qquad
 \eta_m=a_mR\mathbf e_0.
 \label{mem:actualizacion}
\end{equation}
Pour tout \(n\geq1\),
\begin{equation}
 y_n=R^ny_0+\sum_{j=0}^{n-1}R^{n-1-j}\eta_j.
 \label{mem:iteracion}
\end{equation}
En particulier,
\begin{equation}
 y_{12}=y_0+\sum_{j=0}^{11}a_j\mathbf e_j.
 \label{mem:retorno12}
\end{equation}
\end{proposition}
\begin{proof}
Multiplier l'actualisation de \(z_m\) par \(S^{-(m+1)}\) donne
\(y_{m+1}=S^{-1}y_m+a_mS^{-1}\mathbf e_0\), puisque
\(S^{-m}\mathbf e_{m\bmod12}=\mathbf e_0\).
La formule \eqref{mem:iteracion} s'obtient par récurrence.
Pour \(n=12\), \(R^{12}=I\) et
\(R^{11-j}\eta_j=a_jR^{12-j}\mathbf e_0=a_j\mathbf e_j\),
ce qui démontre \eqref{mem:retorno12}.
\end{proof}

Ainsi, douze fenêtres rétablissent le repère et conservent le vecteur des événements.
Le dénombrement ne permet pas à lui seul de retrouver l'ordre de tous les pas
à l'intérieur de chaque fenêtre ; cet ordre demeure dans la trajectoire enregistrée.

Pour préciser le bilan, nous définissons
\[
 D_3=I-S^3,\quad D_4=I-S^4,\quad
 B=D_3^{\mathsf T}D_3+D_4^{\mathsf T}D_4,
\]
\[
 \mathcal M(y)=\tfrac12y^{\mathsf T}By,
 \qquad q(y)=\sum_{j=0}^{11}y_j.
\]
La première est une forme quadratique du registre ; la seconde enregistre
la somme de ses coordonnées. Elles ne sont pas identifiées ici à une énergie ou à une charge physiques.

\begin{proposition}[Bilan produit par l'événement]
\label{mem:prop-balance}
L'actualisation \eqref{mem:actualizacion} satisfait
\begin{equation}
 \begin{aligned}
 \mathcal M(y_{m+1})-\mathcal M(y_m)
   &=a_m(By_m)_0+2a_m^2,\\
 q(y_{m+1})-q(y_m)&=a_m.
 \end{aligned}
 \label{mem:balance}
\end{equation}
\end{proposition}
\begin{proof}
Le développement des produits donne
\(B=4I-S^3-S^{-3}-S^4-S^{-4}\) ; par conséquent,
\(R^{\mathsf T}BR=B\) et \(B_{00}=4\).
Comme \(y_{m+1}=R(y_m+a_m\mathbf e_0)\), la différence quadratique
est \(a_m\mathbf e_0^{\mathsf T}By_m+\tfrac12a_m^2B_{00}\).
La somme des coordonnées est invariante sous la permutation \(R\),
d'où la deuxième égalité.
\end{proof}

\Needspace{20\baselineskip}
\subsubsection{Série génératrice et contrôle uniforme de la mémoire}

La série conserve tous les incréments, et pas uniquement leur somme à la fin
d'un cycle. Avec une variable formelle \(u\), soient
\[
 Y(u)=\sum_{m\geq0}y_mu^m,\qquad
 H_\eta(u)=\sum_{m\geq0}\eta_mu^m.
\]

\begin{proposition}[Résolvante et majoration du reste]
\label{mem:prop-serie}
On a l'identité formelle
\begin{equation}
 Y(u)=(I-uR)^{-1}\bigl(y_0+uH_\eta(u)\bigr).
 \label{mem:serie}
\end{equation}
Les deux séries convergent pour \(|u|<1\). Si \(\ell:\RR^{12}\to\RR\)
est linéaire, \(L=\|\ell\|_{1\to\RR}\), \(M_0=\|y_0\|_1\) et
\(0<r<1\), alors, pour tout entier \(N\geq0\) et \(|u|\leq r\),
\begin{equation}
 \left|\sum_{m>N}\ell(y_m)u^m\right|
 \leq Lr^{N+1}\left(
 \frac{M_0+9(N+1)}{1-r}+\frac{9r}{(1-r)^2}\right).
 \label{mem:cota}
\end{equation}
La série des dérivées converge uniformément dans chaque disque
\(|u|\leq r<1\).
\end{proposition}
\begin{proof}
Multiplier \eqref{mem:actualizacion} par \(u^{m+1}\) et sommer donne
\(Y-y_0=uRY+uH_\eta\). Le terme constant de \(I-uR\) est
inversible, de sorte que son inverse formel est \(\sum_{k\geq0}u^kR^k\).
De plus, \(R\) préserve la norme \(\ell^1\) et
\(\|\eta_m\|_1\leq9\), d'où
\(\|y_m\|_1\leq M_0+9m\).
Pour le reste, on somme les majorants
\(L(M_0+9m)r^m\), en utilisant
\[
 \sum_{m>N}r^m=\frac{r^{N+1}}{1-r},\qquad
 \sum_{m>N}mr^m=r^{N+1}
 \left(\frac{N+1}{1-r}+\frac r{(1-r)^2}\right).
\]
Enfin, \(\sum_{m\geq1}Lm(M_0+9m)r^{m-1}<\infty\)
majore la série dérivée et démontre sa convergence uniforme.
\end{proof}

La variable \(u\) compte les fenêtres de neuf transitions. Son identification
avec la variable d'un autre lecteur exige de conserver cette horloge et l'application
qui relie les deux lectures.

\subsubsection{Élimination de mémoire : opérateur, source et lecteur}

Sur un domaine où l'actualisation est fixée comme \(U:X\to X\),
l'espace vectoriel libre \(V=\QQ^{(X)}\) sur les états admet l'opérateur
\(T\delta_x=\delta_{U(x)}\). L'horloge est incluse dans \(x\) si la
transition en dépend. Cette extension linéaire conserve des états
distincts même lorsqu'ils partagent une phase observable.

Considérons une décomposition \(V=V_0\oplus V_1\), où \(V_1\)
est le secteur auxiliaire que l'on souhaite éliminer, et écrivons
\[
 T=\begin{pmatrix}A&B_1\\C&D\end{pmatrix},\qquad
 v=\binom{v_0}{v_1},\qquad \ell=(\ell_0,\ell_1).
\]
Ici, \(A:V_0\to V_0\), \(B_1:V_1\to V_0\),
\(C:V_0\to V_1\) et \(D:V_1\to V_1\).
Les identités suivantes sont formelles ; elles valent également là où les
résolvantes convergent.

\begin{proposition}[Réduction de Schur avec source et lecteur]
\label{mem:prop-schur}
Soient
\[
 D_u=(I-uD)^{-1},\qquad
 \mathscr F(u)=I-uA-u^2B_1D_uC.
\]
La composante retenue de \((I-uT)^{-1}v\) est
\begin{equation}
 w_0=\mathscr F(u)^{-1}\bigl(v_0+uB_1D_uv_1\bigr),
 \qquad w_1=D_u(v_1+uCw_0).
 \label{mem:schur-fuente}
\end{equation}
Sa lecture complète satisfait
\begin{equation}
 \begin{split}
 \ell(I-uT)^{-1}v
 ={}&(\ell_0+u\ell_1D_uC)
       \mathscr F(u)^{-1}(v_0+uB_1D_uv_1)\\
    &+\ell_1D_uv_1.
 \end{split}
 \label{mem:schur-lector}
\end{equation}
\end{proposition}
\begin{proof}
Les deux équations de \((I-uT)w=v\) sont
\(w_0=v_0+uAw_0+uB_1w_1\) et
\(w_1=v_1+uCw_0+uDw_1\).
Résoudre la seconde et la substituer dans la première donne
\eqref{mem:schur-fuente} ; appliquer \(\ell_0\) et \(\ell_1\)
aux deux composantes donne \eqref{mem:schur-lector}.
Tous les inverses formels existent parce que leurs termes constants
sont des identités.
\end{proof}

Le terme \(u^2B_1D_uC=\sum_{k\geq0}u^{k+2}B_1D^kC\)
décrit une sortie vers la mémoire, \(k\) pas à l'intérieur de celle-ci et
un retour. En revanche, \(u\ell_1D_uC\) lit la mémoire avant ce
retour. Réduire seulement l'opérateur ne conserve donc pas nécessairement
l'observation : la source et le lecteur doivent également être transformés.

La différence est visible dans le cycle déjà construit. Si \(P\) retient
\(\mathbf e_0\), alors \(PRP=0\), mais
\begin{equation}
 \left\langle\mathbf e_0,(I-uR)^{-1}\mathbf e_0\right\rangle
 =\frac1{1-u^{12}},\qquad
 \left\langle\mathbf e_{11},(I-uR)^{-1}\mathbf e_0\right\rangle
 =\frac u{1-u^{12}}.
 \label{mem:ejemplo-ciclo}
\end{equation}
En effet, \(R^n\mathbf e_0=\mathbf e_{-n\bmod12}\) : les
deux séries recueillent respectivement \(n\equiv0\) et
\(n\equiv1\pmod{12}\). La résolvante de la compression est \(I\),
tandis que la compression de la résolvante conserve les retours omis.
Cet exemple correspond au registre de douze fenêtres, et non à la totalité
de l'état du TPK.

\subsubsection{Composition des segments et cohérence de la réduction}

Soient trois transports affines composables
\(F_i(y)=R_i y+b_i\), avec des domaines et codomaines compatibles.
La mémoire produite par le parcours complet est
\begin{equation}
 F_3F_2F_1(y)
 =R_3R_2R_1y+b_3+R_3b_2+R_3R_2b_1.
 \label{mem:cociclo-tres}
\end{equation}
La preuve est la substitution successive des trois formules affines.
Regrouper d'abord les deux premiers segments ou les deux derniers produit
la même expression ; les facteurs transportent chaque incrément depuis
le point où il a pris naissance. Dans \eqref{mem:actualizacion},
\(R_i=R\) et \(b_i\) est l'événement de sa fenêtre.

L'élimination de plusieurs couches conserve une cohérence analogue.
Pour l'expliciter, décomposons maintenant
\(V=V_0\oplus V_1\oplus V_2\) et écrivons
\[
 T=\begin{pmatrix}
 A&B_1&B_2\\C_1&D_{11}&D_{12}\\C_2&D_{21}&D_{22}
 \end{pmatrix},\qquad Q_2=(I-uD_{22})^{-1}.
\]
L'élimination de la troisième composante laisse les quatre blocs
\begin{equation}
 \begin{aligned}
 A'&=A+uB_2Q_2C_2,&
 B_1'&=B_1+uB_2Q_2D_{21},\\
 C_1'&=C_1+uD_{12}Q_2C_2,&
 D_{11}'&=D_{11}+uD_{12}Q_2D_{21}.
 \end{aligned}
 \label{mem:schur-intermedio}
\end{equation}

\begin{proposition}[Transitivité de l'élimination]
Si \(D_*=(D_{ij})_{i,j=1}^2\), le complément obtenu en deux étapes
coïncide avec celui obtenu en éliminant conjointement les deux mémoires :
\begin{equation}
 \begin{split}
 I-uA'-u^2B_1'(I-uD_{11}')^{-1}C_1'
 ={}&I-uA\\
 &-u^2(B_1\ B_2)(I-uD_*)^{-1}\binom{C_1}{C_2}.
 \end{split}
 \label{mem:schur-transitivo}
\end{equation}
\end{proposition}
\begin{proof}
Résoudre la troisième équation de \((I-uT)w=v\) et la substituer dans les
deux premières produit \eqref{mem:schur-intermedio}. Éliminer ensuite la
deuxième composante donne le membre gauche de
\eqref{mem:schur-transitivo}. Résoudre simultanément les deux équations de
mémoire donne le membre droit, par la proposition~\ref{mem:prop-schur}.
La solution formelle est unique parce que \(I-uT\) a pour terme constant
\(I\). Dans les deux procédures, la source et
le lecteur sont également transformés selon \eqref{mem:schur-fuente}--\eqref{mem:schur-lector}.
\end{proof}

Le bilan \eqref{mem:balance} décrit les incréments du registre.
Les identités précédentes permettent de les transporter lors de la réutilisation de la mémoire
dans la clôture dodécaphasique ; elles ne déterminent pas à elles seules des coefficients
analytiques supplémentaires.
